% ============================
% = Verse rendering commands =
% ============================
%
% Structure: Song > Sloka > Sanskrit > Pada > PadaEntry
%                         > Translation (ftrans)
%                         > Footnotes

\usepackage{glossy}

\definecolor{glossgray}{HTML}{4D4D4D}
\gsSetFirstStyle{\Large\bfseries}
\gsSetSecondStyle{\color{glossgray}\itshape\normalsize}
\newlength\gsFirstAscent
\newlength\gsSecondAscent
\AtBeginDocument{%
	\settoheight\gsFirstAscent{\Large M}%
	\setlength\gsFirstHeight{\gsFirstAscent}%
	\settoheight\gsSecondAscent{\normalsize M}%
	\setlength\gsSecondHeight{\gsSecondAscent}%
}

% -- Song: a named section (e.g. Gurvaṣṭakam) --
%   Stores the song title for display; running headers use \leftmark
%   (chapter title) and \booktitle (book title) instead.
\newenvironment{song}[1]
{\def\songtitle{#1}}
{}

% Defensive default so the sloka environment never hits an undefined
% \songtitle if a sloka is ever typeset outside a song.
\def\songtitle{}

% -- Song rule: the chapter's swelled rule shrunk to a hairline, drawn
%   in the gloss gray --
\newcommand{\songrule}{%
	\tikz \filldraw[scale=.015,thin,draw=glossgray,fill=glossgray]%
	(0,0) -- (100,0.7) -- (200,0.7) -- (300,0) --%
	(200,-0.7) -- (100,-0.7) -- cycle;%
}

% -- Bengali-song marker: a small gray superscript "B" appended to the
%   title of songs whose text is in Bengali (language: bn in the verse
%   YAML). Set in gloss gray to match the author line and running head.
\newcommand{\bengalimark}{\textsuperscript{\textcolor{glossgray}{B}}}

% -- No-translation rule: like \songrule but slightly narrower, marks a
%   sloka rendered without a translation (e.g. gaura-arati refrains) --
\newcommand{\notransrule}{%
	\tikz \filldraw[scale=.015,thin,draw=glossgray,fill=glossgray]%
	(0,0) -- (85,0.7) -- (170,0.7) -- (255,0) --%
	(170,-0.7) -- (85,-0.7) -- cycle;%
}

% -- Song header: title, author, hairline rule --
%   Emitted inside the first sloka's unbreakable box, so it can never be
%   separated from the verse (the verse box supplies the leading space).
%   Blank author suppresses the author line. \nointerlineskip on both sides
%   of the rule (prevdepth = -1000pt) removes the ~1ex natural line gap, so
%   exactly 1ex separates the rule from the text above and the verse
%   below.
\ExplSyntaxOn
\newcommand{\songheader}[2]{%
	{\centering\LARGE\bfseries #1\par}%
	\tl_if_blank:nF {#2}
	{%
		\vspace{0.5ex}%
		{\centering\Large\color{glossgray}#2\par}%
	}%
	\nointerlineskip
	\vspace{1ex}%
	{\centering\songrule\par}%
	\nointerlineskip
	\vspace{1ex}%
}
\ExplSyntaxOff

% -- Sloka: a numbered verse, kept whole --
%   The whole verse (number, sanskrit, translation) is set in an
%   unbreakable box, so it can never split across a page. The blank
%   argument → no header; any other → vertical skip + header.
%
%   Running headers use \leftmark (chapter title) and \booktitle
%   (book title), so no song mark is needed here.
\ExplSyntaxOn
% The verse box is built in a register so its real size is known before it
% is committed to the page.
\newbox\songbookverse
\newenvironment{sloka}[1]
{
	\par
	\vspace{1ex} % space inbetween the verse boxes
	\setbox\songbookverse=\vbox\bgroup
	\vspace{1ex} % space inside the verse box
	\tl_if_blank:oTF {#1}
	{}
	{
		{\centering\Large(#1)\par}
		\vspace{1ex}
	}
}
{
	\egroup
	\box\songbookverse
}
\ExplSyntaxOff

% -- Sanskrit: semantic wrapper (currently no-op) --
\newenvironment{sanskrit}{}{}

% -- Pada: a line of word-by-word glosses --
\newenvironment{pada}%
{\begin{glossy}\begin{nscenter}}%
{\end{nscenter}\end{glossy}}%

\newcommand{\padaEntry}[2]{\gloss{#1}{#2}}


% -- Translation: full-width ragged-right text --
\newcommand{\ftransstyle}{%
	\raggedright
	\hyphenpenalty=5000
}
\newcommand{\ftrans}[1]{%
	\vspace{0.9ex}
	{\ftransstyle\normalsize #1\par}
}

% -- Footnote: term + em-dash + note with a hanging indent so wrapped
%   lines are indented and each note reads as one block --
\newcommand{\footnoteEntry}[2]{%
	\par\noindent
	\hangindent=2em\hangafter=1
	\textsc{#1}\enspace\textemdash\enspace#2\par
}
